\documentclass[11pt,a4paper]{article}
\usepackage{fontspec,amsmath,amssymb,amsthm,microtype}
\setmainfont{lmroman10-regular.otf}[BoldFont=lmroman10-bold.otf,ItalicFont=lmroman10-italic.otf,BoldItalicFont=lmroman10-bolditalic.otf]
\usepackage{xcolor}
\usepackage[most]{tcolorbox}
\usepackage{mathrsfs}
\providecommand{\lt}{<}
\providecommand{\gt}{>}
\setmonofont{DejaVuSansMono.ttf}[BoldFont=DejaVuSansMono-Bold.ttf,ItalicFont=DejaVuSansMono-Oblique.ttf,BoldItalicFont=DejaVuSansMono-BoldOblique.ttf,Scale=0.86]
\newfontfamily\leanfallback{STIXTwoMath-Regular.otf}[Scale=0.86]
\definecolor{leanrule}{gray}{0.62}
\definecolor{leanback}{gray}{0.962}
\definecolor{leantitle}{gray}{0.22}
\definecolor{leanmuted}{gray}{0.38}
\newtcolorbox{leanbox}[1]{enhanced,breakable,colback=leanback,colframe=leanrule,
  boxrule=0.4pt,arc=0.6mm,left=2.4mm,right=2.4mm,top=1.8mm,bottom=1.8mm,
  before skip=0.8em,after skip=1em,fonttitle=\small\bfseries,coltitle=leantitle,
  colbacktitle=leanback,titlerule=0.3pt,title={#1}}
\newlength{\leancharwidth}
\newenvironment{leancode}{\par\smallskip\begingroup\ttfamily\footnotesize\raggedright
  \settowidth{\leancharwidth}{M}\setlength{\parskip}{0pt}\setlength{\parindent}{0pt}
  \frenchspacing}{\par\endgroup\smallskip}
\newcommand{\leanline}[2]{\hangindent=\dimexpr#1\leancharwidth+4\leancharwidth\relax
  \hangafter=1\noindent\hspace*{#1\leancharwidth}#2\par}
\newcommand{\leansource}[1]{\par{\footnotesize\color{leanmuted}#1}\par}
\newcommand{\leanlabel}[1]{\par\smallskip\noindent{\small\bfseries #1}\par}
\usepackage[margin=25mm]{geometry}
\usepackage[unicode,hidelinks]{hyperref}
\hypersetup{pdfauthor={LangevinLean Project}}
\setlength{\emergencystretch}{3em}
\setlength{\parskip}{0.4em}
\setlength{\parindent}{0pt}
\allowdisplaybreaks[2]
\raggedbottom
\title{\texorpdfstring{Time differentiation and local space-time integrability of the Gaussian mixture}{Time differentiation and local space-time integrability of the Gaussian mixture}}
\author{LangevinLean Project}
\date{}
\begin{document}\maketitle
Draft mathematical exposition; these notes may contain mistakes.


Let \(g_x=\nabla U(x)\), \(c_t(x)=x-tg_x\), and \(K_t(x,y)=(4\pi t)^{-d/2}\exp(-\|y-c_t(x)\|^2/(4t))\). At \(t>0\), differentiation of the normalizer and exponent gives \(\partial_tK_t=K_t[-d/(2t)+\|y-x\|^2/(4t^2)-\|g_x\|^2/4]\). No derivative of \(g_x\) is taken.



Assume \(G_0=\int\|g_x\|^2d\mu<\infty\). On a neighborhood of each positive time, a bound by a constant multiple of \(1+\|g_x\|^2\) justifies differentiation under the initial-law integral and proves \(\partial_tp_t(y)=\int\partial_tK_t(x,y)d\mu(x)\). Retaining the Gaussian factor instead gives a bound integrable in both starting point and observation. Its explicit integral proves local space-time integrability of \(\partial_tp_t\) on intervals bounded away from zero. The entropy chain rule uses these estimates after initial gradient energy has been deduced from finite initial relative entropy.


\section{\texorpdfstring{The Gaussian normalizer as an explicit time power}{The Gaussian normalizer as an explicit time power}}


The proposal normalizer is defined as the reciprocal of the Gaussian integral. Gaussian Fourier integration evaluates \(\int e^{-a\|z\|^2}dz=(\pi/a)^{d/2}\) for \(a>0\). Substitute \(a=1/(4t)\), simplify \(\pi/a=4\pi t\), and invert the real power to obtain \(N_t=(4\pi t)^{-d/2}\). Positivity of \(t\) justifies the Gaussian formula and the real-power inversion.



The Gaussian normalizer has time dependence \((4\pi t)^{-d/2}\). Expressing it as a real power separates the positive constant from the time power and permits ordinary scalar differentiation.

\begin{leanbox}{Lean: \texttt{gaussianProposalNormalizer\_\allowbreak{}eq\_\allowbreak{}rpow}}
\leansource{\texttt{gaussianProposalNormalizer\_\allowbreak{}eq\_\allowbreak{}rpow} \quad UnadjustedLangevin/\allowbreak{}LSI/\allowbreak{}GaussianTimeDerivative.lean, lines 26–34}
\begin{leancode}
\leanline{0}{theorem\ gaussianProposalNormalizer\_eq\_rpow\ (t\ :\ ℝ)\ (ht\ :\ 0\ <\ t)\ :}
\leanline{4}{ULA.gaussianProposalNormalizer\ (H\ :=\ E)\ t\ =}
\leanline{6}{(4\ *\ Real.pi\ *\ t)\ \textasciicircum{}\ (-(Module.finrank\ ℝ\ E\ :\ ℝ)\ /\ 2)}
\end{leancode}
\leanlabel{How the Lean proof follows the argument}
\texttt{gaussian\allowbreak{}Proposal\allowbreak{}Normalizer\_\allowbreak{}eq\_\allowbreak{}rpow}: the intermediate fact is \texttt{heq}, in source order; the source steps apply or rewrite \texttt{Gaussian\allowbreak{}Fourier.\allowbreak{}integral\_\allowbreak{}rexp\_\allowbreak{}neg\_\allowbreak{}mul\_\allowbreak{}sq\_\allowbreak{}norm}, \texttt{Module.\allowbreak{}finrank}. The remaining scalar rearrangements are checked by \texttt{field\_\allowbreak{}simp}, \texttt{ring}.
\end{leanbox}


Earlier results used here: \cite[Section~3.37.1]{LangevinLean2026ULASC}; \cite[Section~3.37.12]{LangevinLean2026ULASC}; \cite[Section~3.37.2]{LangevinLean2026ULASC}.

\section{\texorpdfstring{The time derivative of the Gaussian normalizer}{The time derivative of the Gaussian normalizer}}


Fix \(t>0\). On a neighborhood of \(t\), every time \(s\) remains positive, and the Gaussian integral formula gives \(N_s=(4\pi s)^{-d/2}\). The real-power chain rule therefore gives

\[N_t'=(-d/2)(4\pi t)^{-d/2-1}(4\pi).\]

Because \(4\pi t>0\), the identity \((4\pi t)^{p-1}=(4\pi t)^p/(4\pi t)\) is valid. Substituting \(p=-d/2\) simplifies the derivative to

\[N_t'=-\frac d{2t}N_t.\]

The local equality with the explicit power transfers this derivative to the normalizer defined by the Gaussian integral. Thus the density normalization contributes logarithmic derivative \(-d/(2t)\).

\begin{leanbox}{Lean: \texttt{hasDerivAt\_\allowbreak{}gaussianProposalNormalizer}}
\leansource{\texttt{hasDerivAt\_\allowbreak{}gaussianProposalNormalizer} \quad UnadjustedLangevin/\allowbreak{}LSI/\allowbreak{}GaussianTimeDerivative.lean, lines 39–64}
\begin{leancode}
\leanline{0}{theorem\ hasDerivAt\_gaussianProposalNormalizer\ (t\ :\ ℝ)\ (ht\ :\ 0\ <\ t)\ :}
\leanline{4}{HasDerivAt\ (ULA.gaussianProposalNormalizer\ (H\ :=\ E))}
\leanline{6}{(-(Module.finrank\ ℝ\ E\ :\ ℝ)\ /\ (2\ *\ t)\ *\ ULA.gaussianProposalNormalizer\ (H\ :=\ E)\ t)\ t}
\end{leancode}
\leanlabel{How the Lean proof follows the argument}
\texttt{has\allowbreak{}Deriv\allowbreak{}At\_\allowbreak{}gaussian\allowbreak{}Proposal\allowbreak{}Normalizer}: the intermediate facts include \texttt{p}, \texttt{hp}, \texttt{hd}, \texttt{he}, \texttt{hder}, in source order; the source steps apply or rewrite \texttt{gaussian\allowbreak{}Proposal\allowbreak{}Normalizer\_\allowbreak{}eq\_\allowbreak{}rpow}, \texttt{Real.\allowbreak{}rpow\_\allowbreak{}sub}, \texttt{id\_\allowbreak{}eq}, \texttt{mul\_\allowbreak{}one}, \texttt{hd'.\allowbreak{}congr\_\allowbreak{}of\_\allowbreak{}eventually\allowbreak{}Eq}. The remaining scalar rearrangements are checked by \texttt{field\_\allowbreak{}simp}, \texttt{ring}.
\end{leanbox}


Earlier results used here: \cite[Section~3.37.12]{LangevinLean2026ULASC}.

\section{\texorpdfstring{The frozen-drift Gaussian kernel}{The frozen-drift Gaussian kernel}}


The real component kernel is \(K_t(x,y)=N_t\exp(-\|y-x+tg_x\|^2/(4t))\). The gradient \(g_x\) is evaluated at the initial state \(x\), and is constant when differentiating in \(t\) or \(y\). Its integral over \(y\) is one at positive time by the previously proved Gaussian proposal normalization.



The frozen kernel is the Gaussian density conditional on the initial state. Its center moves linearly with time through the fixed vector \(\nabla U(x)\), while its covariance grows as \(2tI_d\).

\begin{leanbox}{Lean: \texttt{frozenGaussianKernel}}
\leansource{\texttt{frozenGaussianKernel} \quad UnadjustedLangevin/\allowbreak{}LSI/\allowbreak{}GaussianTimeDerivative.lean, lines 67–69}
\begin{leancode}
\leanline{0}{def\ frozenGaussianKernel\ (U\ :\ E\ →\ ℝ)\ (t\ :\ ℝ)\ (x\ y\ :\ E)\ :\ ℝ\ :=}
\leanline{2}{ULA.gaussianProposalNormalizer\ (H\ :=\ E)\ t\ *}
\leanline{4}{Real.exp\ (-‖y\ -\ (x\ -\ t\ •\ gradient\ U\ x)‖\ \textasciicircum{}\ 2\ /\ (4\ *\ t))}
\end{leancode}
\leanlabel{How the Lean proof follows the argument}
\texttt{frozen\allowbreak{}Gaussian\allowbreak{}Kernel} is a definition with defining expression \texttt{E) t *}. The right-hand side specifies the object used in the mathematical statement.
\end{leanbox}


Earlier results used here: \cite[Section~3.37.12]{LangevinLean2026ULASC}.

\section{\texorpdfstring{Its explicit candidate time derivative}{Its explicit candidate time derivative}}


Define the explicit candidate derivative \(\dot K_t=K_t[-d/(2t)+\|y-x\|^2/(4t^2)-\|g_x\|^2/4]\). The dimension term comes from the normalizer; the remaining terms come from differentiating the exponent. In the expanded exponent, the mixed inner-product term is independent of time, which explains its absence from this expression. The following theorem verifies that this formula is the derivative.



The derivative field is defined by the explicit normalizer and exponent formula. Naming this function before differentiation makes later measurability and integrability statements refer to one fixed spacetime integrand.

\begin{leanbox}{Lean: \texttt{frozenGaussianKernelDot}}
\leansource{\texttt{frozenGaussianKernelDot} \quad UnadjustedLangevin/\allowbreak{}LSI/\allowbreak{}GaussianTimeDerivative.lean, lines 72–75}
\begin{leancode}
\leanline{0}{def\ frozenGaussianKernelDot\ (U\ :\ E\ →\ ℝ)\ (t\ :\ ℝ)\ (x\ y\ :\ E)\ :\ ℝ\ :=}
\leanline{2}{frozenGaussianKernel\ U\ t\ x\ y\ *}
\leanline{4}{(-(Module.finrank\ ℝ\ E\ :\ ℝ)\ /\ (2\ *\ t)\ +\ ‖y\ -\ x‖\ \textasciicircum{}\ 2\ /\ (4\ *\ t\ \textasciicircum{}\ 2)\ -}
\leanline{6}{‖gradient\ U\ x‖\ \textasciicircum{}\ 2\ /\ 4)}
\end{leancode}
\leanlabel{How the Lean proof follows the argument}
\texttt{frozen\allowbreak{}Gaussian\allowbreak{}Kernel\allowbreak{}Dot} is a definition with defining expression \texttt{frozen\allowbreak{}Gaussian\allowbreak{}Kernel U t x y *}. The right-hand side specifies the object used in the mathematical statement.
\end{leanbox}

\section{\texorpdfstring{Expanding the frozen Gaussian exponent}{Expanding the frozen Gaussian exponent}}


Set \(z=y-x\). The squared-norm identity gives \(\|z+sg_x\|^2=\|z\|^2+2s\langle z,g_x\rangle+s^2\|g_x\|^2\). For \(s\ne0\), division by \(-4s\) yields \(-\|z+sg_x\|^2/(4s)=-\|z\|^2/(4s)-\langle z,g_x\rangle/2-s\|g_x\|^2/4\). The vector identity and scalar denominator cancellation are proved separately in the proof calculation.



Expanding \(\|y-x+t g\|^2\) and dividing by \(4t\) yields \(\|y-x\|^2/(4t)+\langle y-x,g\rangle/2+t\|g\|^2/4\). The middle term has zero time derivative, simplifying the formula.

\begin{leanbox}{Lean: \texttt{frozen\_\allowbreak{}exponent\_\allowbreak{}eq}}
\leansource{\texttt{frozen\_\allowbreak{}exponent\_\allowbreak{}eq} \quad UnadjustedLangevin/\allowbreak{}LSI/\allowbreak{}GaussianTimeDerivative.lean, lines 78–86}
\begin{leancode}
\leanline{0}{private\ theorem\ frozen\_exponent\_eq\ (U\ :\ E\ →\ ℝ)\ (s\ :\ ℝ)\ (hs\ :\ s\ ≠\ 0)\ (x\ y\ :\ E)\ :}
\leanline{4}{-‖y\ -\ (x\ -\ s\ •\ gradient\ U\ x)‖\ \textasciicircum{}\ 2\ /\ (4\ *\ s)\ =}
\leanline{6}{-(‖y\ -\ x‖\ \textasciicircum{}\ 2\ /\ 4)\ /\ s\ -\ inner\ ℝ\ (y\ -\ x)\ (gradient\ U\ x)\ /\ 2\ -}
\leanline{8}{s\ *\ ‖gradient\ U\ x‖\ \textasciicircum{}\ 2\ /\ 4}
\end{leancode}
\leanlabel{How the Lean proof follows the argument}
\texttt{frozen\_\allowbreak{}exponent\_\allowbreak{}eq}: the intermediate fact is \texttt{he}, in source order; the source steps apply or rewrite \texttt{norm\_\allowbreak{}add\_\allowbreak{}sq\_\allowbreak{}real}, \texttt{inner\_\allowbreak{}smul\_\allowbreak{}right}, \texttt{Real.\allowbreak{}norm\_\allowbreak{}eq\_\allowbreak{}abs}, \texttt{mul\_\allowbreak{}pow}, \texttt{sq\_\allowbreak{}abs}. The remaining scalar rearrangements are checked by \texttt{field\_\allowbreak{}simp}, \texttt{ring}.
\end{leanbox}

\section{\texorpdfstring{Differentiating the frozen Gaussian kernel}{Differentiating the frozen Gaussian kernel}}


Write the expanded exponent as \(-A/s-B-sC\), where \(A=\|y-x\|^2/4\), \(B=\langle y-x,g_x\rangle/2\), and \(C=\|g_x\|^2/4\). Its derivative at positive \(t\) is \(A/t^2-C\). The expansion agrees with the original exponent on a neighborhood avoiding zero, so this is also the derivative of the original expression.



Apply the exponential chain rule and multiply by the normalizer. Its derivative contributes \(-dK_t/(2t)\); the exponential derivative contributes \(K_t(A/t^2-C)\). Collecting terms gives exactly its explicit time derivative. No derivative of \(g_x\), and therefore no Hessian, is taken.



Combine the normalizer derivative with the derivative of the expanded exponent. The cross term contributes no time derivative, leaving exactly the stated squared-position and squared-gradient terms.

\begin{leanbox}{Lean: \texttt{hasDerivAt\_\allowbreak{}frozenGaussianKernel}}
\leansource{\texttt{hasDerivAt\_\allowbreak{}frozenGaussianKernel} \quad UnadjustedLangevin/\allowbreak{}LSI/\allowbreak{}GaussianTimeDerivative.lean, lines 90–114}
\begin{leancode}
\leanline{0}{theorem\ hasDerivAt\_frozenGaussianKernel\ (U\ :\ E\ →\ ℝ)\ (t\ :\ ℝ)\ (ht\ :\ 0\ <\ t)\ (x\ y\ :\ E)\ :}
\leanline{4}{HasDerivAt\ (fun\ s\ ↦\ frozenGaussianKernel\ U\ s\ x\ y)}
\leanline{6}{(frozenGaussianKernelDot\ U\ t\ x\ y)\ t}
\end{leancode}
\leanlabel{How the Lean proof follows the argument}
\texttt{has\allowbreak{}Deriv\allowbreak{}At\_\allowbreak{}frozen\allowbreak{}Gaussian\allowbreak{}Kernel}: the intermediate facts include \texttt{A}, \texttt{B}, \texttt{C}, \texttt{heq}, \texttt{hexp}, in source order; the source steps apply or rewrite \texttt{div\_\allowbreak{}eq\_\allowbreak{}mul\_\allowbreak{}inv}, \texttt{Pi.\allowbreak{}sub\_\allowbreak{}apply}, \texttt{mul\_\allowbreak{}assoc}, \texttt{frozen\_\allowbreak{}exponent\_\allowbreak{}eq}, \texttt{ht.\allowbreak{}ne'}. The remaining scalar rearrangements are checked by \texttt{ring}.
\end{leanbox}

\section{\texorpdfstring{An absolute bound on the kernel time derivative}{An absolute bound on the kernel time derivative}}


Put \(r=y-x+tg_x\) and \(e=e^{-\|r\|^2/(4t)}\). The triangle inequality and \((a+b)^2\le2a^2+2b^2\) give \(\|y-x\|^2\le2\|r\|^2+2t^2\|g_x\|^2\). The Gaussian exponential bound gives \(\|r\|^2e\le4t\), while \(e\le1\). Consequently \(\|y-x\|^2e\le8t+2t^2\|g_x\|^2\).



Bound the absolute value of each derivative-coefficient term by their nonnegative sum, multiply by \(e\), and insert these bounds. The resulting upper bound is \(|\dot K_t(x,y)|\le N_t[(d+4)/t+\|g_x\|^2]\). It is uniform in \(y\), and is integrable in the initial variable whenever \(G_0<\infty\). Its purpose is time differentiation under the initial-state integral, and will be replaced by a Gaussian-weighted bound for output-space integration.



Take absolute values term by term in the explicit derivative. Nonnegativity of the kernel allows its factor to be retained, producing a sum of Gaussian-weighted nonnegative energy terms.

\begin{leanbox}{Lean: \texttt{norm\_\allowbreak{}frozenGaussianKernelDot\_\allowbreak{}le}}
\leansource{\texttt{norm\_\allowbreak{}frozenGaussianKernelDot\_\allowbreak{}le} \quad UnadjustedLangevin/\allowbreak{}LSI/\allowbreak{}GaussianTimeDerivative.lean, lines 119–174}
\begin{leancode}
\leanline{0}{theorem\ norm\_frozenGaussianKernelDot\_le\ (U\ :\ E\ →\ ℝ)\ (t\ :\ ℝ)\ (ht\ :\ 0\ <\ t)\ (x\ y\ :\ E)\ :}
\leanline{4}{‖frozenGaussianKernelDot\ U\ t\ x\ y‖\ ≤}
\leanline{6}{ULA.gaussianProposalNormalizer\ (H\ :=\ E)\ t\ *}
\leanline{8}{(((Module.finrank\ ℝ\ E\ :\ ℝ)\ +\ 4)\ /\ t\ +\ ‖gradient\ U\ x‖\ \textasciicircum{}\ 2)}
\end{leancode}
\leanlabel{How the Lean proof follows the argument}
\texttt{norm\_\allowbreak{}frozen\allowbreak{}Gaussian\allowbreak{}Kernel\allowbreak{}Dot\_\allowbreak{}le}: the intermediate facts include \texttt{r}, \texttt{e}, \texttt{A}, \texttt{ht2}, \texttt{hformula}, in source order; the source steps apply or rewrite \texttt{div\_\allowbreak{}nonpos\_\allowbreak{}of\_\allowbreak{}nonpos\_\allowbreak{}of\_\allowbreak{}nonneg}, \texttt{neg\_\allowbreak{}nonpos.\allowbreak{}mpr}, \texttt{abs\_\allowbreak{}mul}, \texttt{mul\_\allowbreak{}le\_\allowbreak{}mul\_\allowbreak{}of\_\allowbreak{}nonneg\_\allowbreak{}left}, \texttt{ULA.\allowbreak{}gaussian\allowbreak{}Proposal\allowbreak{}Normalizer\_\allowbreak{}pos}. The remaining scalar rearrangements are checked by \texttt{field\_\allowbreak{}simp}, \texttt{ring}, \texttt{linarith}, \texttt{nlinarith}.
\end{leanbox}


Earlier results used here: \cite[Section~3.37.12]{LangevinLean2026ULASC}; \cite[Section~3.37.14]{LangevinLean2026ULASC}.

\section{\texorpdfstring{Measurability of the normalizer}{Measurability of the normalizer}}


At positive time the Gaussian normalizer has the continuous explicit form \((4\pi t)^{-d/2}\), hence is measurable. Assigning an arbitrary value at zero gives a measurable extension if a whole-time parameter field is desired. Only the positive-time values are used in the Gaussian density and derivative identities.



At positive time the normalizer is continuous, hence measurable. Its reciprocal-power expression also permits a measurable extension when the proof organizes all times in one ambient domain.

\begin{leanbox}{Lean: \texttt{gaussianProposalNormalizer\_\allowbreak{}measurable}}
\leansource{\texttt{gaussianProposalNormalizer\_\allowbreak{}measurable} \quad UnadjustedLangevin/\allowbreak{}LSI/\allowbreak{}GaussianTimeDerivative.lean, lines 179–183}
\begin{leancode}
\leanline{0}{theorem\ gaussianProposalNormalizer\_measurable\ :}
\leanline{4}{Measurable\ (ULA.gaussianProposalNormalizer\ (H\ :=\ E))}
\end{leancode}
\leanlabel{How the Lean proof follows the argument}
\texttt{gaussian\allowbreak{}Proposal\allowbreak{}Normalizer\_\allowbreak{}measurable}: the intermediate fact is \texttt{hm}, in source order; the source steps apply or rewrite \texttt{hm.\allowbreak{}integral\_\allowbreak{}prod\_\allowbreak{}right.\allowbreak{}measurable.\allowbreak{}inv}.
\end{leanbox}


Earlier results used here: \cite[Section~3.37.12]{LangevinLean2026ULASC}.

\section{\texorpdfstring{Joint measurability of the kernel time derivative}{Joint measurability of the kernel time derivative}}


Lipschitzness makes \(x\mapsto g_x\) continuous and measurable. Combine this with the measurable normalizer and the explicit algebraic formula for \(\dot K_t\). Products, norms, scalar inverses, and exponentials preserve measurability, giving joint measurability in \((t,y,x)\). This is the input needed to integrate out the initial variable while retaining measurable time-space dependence.



The explicit derivative uses only measurable gradient, norms, scalar operations, and the Gaussian exponential. Their composition is jointly measurable in time, initial state, and observation.

\begin{leanbox}{Lean: \texttt{frozenGaussianKernelDot\_\allowbreak{}measurable}}
\leansource{\texttt{frozenGaussianKernelDot\_\allowbreak{}measurable} \quad UnadjustedLangevin/\allowbreak{}LSI/\allowbreak{}GaussianTimeDerivative.lean, lines 186–192}
\begin{leancode}
\leanline{0}{theorem\ frozenGaussianKernelDot\_measurable\ (U\ :\ E\ →\ ℝ)\ (L\ :\ ℝ≥0)}
\leanline{4}{(hLip\ :\ LipschitzWith\ L\ (gradient\ U))\ :}
\leanline{4}{Measurable\ (fun\ p\ :\ (ℝ\ ×\ E)\ ×\ E\ ↦\ frozenGaussianKernelDot\ U\ p.1.1\ p.2\ p.1.2)}
\end{leancode}
\leanlabel{How the Lean proof follows the argument}
\texttt{frozen\allowbreak{}Gaussian\allowbreak{}Kernel\allowbreak{}Dot\_\allowbreak{}measurable}: the intermediate facts include \texttt{hg}, \texttt{hn}, in source order.
\end{leanbox}

\section{\texorpdfstring{The candidate time derivative of the mixture}{The candidate time derivative of the mixture}}


Define \(\dot p_t(y)=\int\dot K_t(x,y)\,d\mu(x)\). At this point it is an ordinary parameter integral. The subsequent differentiation theorem proves it is the derivative of \(p_t(y)\), and the energy estimates prove its integrability, so differentiation and integrability remain conclusions to prove.



Integrating the explicit kernel derivative over the initial probability defines the candidate mixture time derivative. Its later identification as a derivative requires the uniform integrable bound proved separately.

\begin{leanbox}{Lean: \texttt{interpolationLebesgueDensityTimeDerivative}}
\leansource{\texttt{interpolationLebesgueDensityTimeDerivative} \quad UnadjustedLangevin/\allowbreak{}LSI/\allowbreak{}GaussianTimeDerivative.lean, lines 195–197}
\begin{leancode}
\leanline{0}{def\ interpolationLebesgueDensityTimeDerivative}
\leanline{4}{(U\ :\ E\ →\ ℝ)\ (μ\ :\ Measure\ E)\ (t\ :\ ℝ)\ (y\ :\ E)\ :\ ℝ\ :=}
\leanline{2}{∫\ x,\ frozenGaussianKernelDot\ U\ t\ x\ y\ ∂μ}
\end{leancode}
\leanlabel{How the Lean proof follows the argument}
\texttt{interpolation\allowbreak{}Lebesgue\allowbreak{}Density\allowbreak{}Time\allowbreak{}Derivative} is a definition with defining expression \texttt{\ensuremath{\int} x, frozen\allowbreak{}Gaussian\allowbreak{}Kernel\allowbreak{}Dot U t x y \ensuremath{\partial}\ensuremath{\mu}}. The right-hand side specifies the object used in the mathematical statement.
\end{leanbox}

\section{\texorpdfstring{Joint measurability of the mixture derivative}{Joint measurability of the mixture derivative}}


The component derivative is jointly measurable in \((t,y,x)\). Its strongly measurable parameter integral over the probability measure \(\mu\) is therefore measurable in \((t,y)\). This gives a single measurable version of \(\dot p\) for spacetime Fubini and the entropy chain rule.



Parameter integration of the jointly measurable derivative field gives a measurable candidate in time and observation. Absolute integrability, when invoked later, ensures it is also the ordinary integral required by the chain rule.

\begin{leanbox}{Lean: \texttt{interpolationLebesgueDensityTimeDerivative\_\allowbreak{}measurable}}
\leansource{\texttt{interpolationLebesgueDensityTimeDerivative\_\allowbreak{}measurable} \quad UnadjustedLangevin/\allowbreak{}LSI/\allowbreak{}GaussianTimeDerivative.lean, lines 200–205}
\begin{leancode}
\leanline{0}{theorem\ interpolationLebesgueDensityTimeDerivative\_measurable}
\leanline{4}{(U\ :\ E\ →\ ℝ)\ (μ\ :\ Measure\ E)\ [IsProbabilityMeasure\ μ]}
\leanline{4}{(L\ :\ ℝ≥0)\ (hLip\ :\ LipschitzWith\ L\ (gradient\ U))\ :}
\leanline{4}{Measurable\ (fun\ p\ :\ ℝ\ ×\ E\ ↦\ interpolationLebesgueDensityTimeDerivative\ U\ μ\ p.1\ p.2)}
\end{leancode}
\leanlabel{How the Lean proof follows the argument}
\texttt{interpolation\allowbreak{}Lebesgue\allowbreak{}Density\allowbreak{}Time\allowbreak{}Derivative\_\allowbreak{}measurable}: the source steps apply or rewrite \texttt{frozen\allowbreak{}Gaussian\allowbreak{}Kernel\allowbreak{}Dot\_\allowbreak{}measurable}.
\end{leanbox}

\section{\texorpdfstring{The time derivative of the mixture}{The time derivative of the mixture}}


Fix \(t>0\). Continuity of \(N_s\) gives a neighborhood \(S\) of \(t\) on which \(s>t/2\) and \(N_s<N_t+1\). The derivative bound implies, for every \(s\in S\) and fixed observation \(y\),

\[|\dot K_s(x,y)|\le(N_t+1)\left[\frac{2(d+4)}t+\|g_x\|^2\right].\]

The right side is integrable in \(x\) because \(G_0<\infty\) and \(\mu\) is a probability. The component kernel at \(t\) is also integrable, since it is bounded by \(N_t\). Each component has the time derivative already computed.



Differentiation under an integral with this common integrable derivative bound gives both integrability of \(\dot K_t(\cdot,y)\) and

\[\frac{d}{dt}\int K_t(x,y)d\mu(x)=\int\dot K_t(x,y)d\mu(x).\]

The integral on the left agrees with the mixture density on the positive-time neighborhood. Hence \(\partial_tp_t(y)=\dot p_t(y)\) at every observation. This argument does not require invertibility of the Euler map or a step-size restriction.

\begin{leanbox}{Lean: \texttt{hasDerivAt\_\allowbreak{}interpolationLebesgueDensity\_\allowbreak{}time}}
\leansource{\texttt{hasDerivAt\_\allowbreak{}interpolationLebesgueDensity\_\allowbreak{}time} \quad UnadjustedLangevin/\allowbreak{}LSI/\allowbreak{}GaussianTimeDerivative.lean, lines 210–275}
\begin{leancode}
\leanline{0}{theorem\ hasDerivAt\_interpolationLebesgueDensity\_time}
\leanline{4}{(U\ :\ E\ →\ ℝ)\ (μ\ :\ Measure\ E)\ [IsProbabilityMeasure\ μ]}
\leanline{4}{(L\ :\ ℝ≥0)\ (hLip\ :\ LipschitzWith\ L\ (gradient\ U))}
\leanline{4}{(hG\ :\ Integrable\ (fun\ x\ ↦\ ‖gradient\ U\ x‖\ \textasciicircum{}\ 2)\ μ)}
\leanline{4}{(t\ :\ ℝ)\ (ht\ :\ 0\ <\ t)\ (y\ :\ E)\ :}
\leanline{4}{Integrable\ (fun\ x\ ↦\ frozenGaussianKernelDot\ U\ t\ x\ y)\ μ\ ∧}
\leanline{6}{HasDerivAt\ (fun\ s\ ↦\ interpolationLebesgueDensity\ U\ μ\ s\ y)}
\leanline{8}{(interpolationLebesgueDensityTimeDerivative\ U\ μ\ t\ y)\ t}
\end{leancode}
\leanlabel{How the Lean proof follows the argument}
\texttt{interpolation\allowbreak{}Lebesgue\_\allowbreak{}fisher\_\allowbreak{}bound} uses the differentiated mixture formula and weighted Cauchy–Schwarz to bound the mixture score by the averaged component-score energy. \texttt{has\allowbreak{}Deriv\allowbreak{}At\_\allowbreak{}interpolation\allowbreak{}Lebesgue\allowbreak{}Density\_\allowbreak{}time} justifies differentiation of the Gaussian probability mixture with the positive-time kernel domination estimates. \texttt{interpolation\allowbreak{}Lebesgue\allowbreak{}Density\_\allowbreak{}weak\_\allowbreak{}identity} applies the Gaussian flux integration-by-parts result under its bounded-test and finite-energy premises. The statements use the same interpolation density and initial gradient moment, so these results can be combined with the relative-density identity and the log-Sobolev inequality. Their conclusions concern positive times; none assumes differentiability at zero.
\end{leanbox}


Earlier results used here: \cite[Section~3.37.12]{LangevinLean2026ULASC}; \cite[Section~3.37.14]{LangevinLean2026ULASC}.

\section{\texorpdfstring{Equality of the two Gaussian kernel formulas}{Equality of the two Gaussian kernel formulas}}


The spatial Fisher calculation uses precision notation \(q_a(c,y)=Z_a^{-1}e^{-a\|y-c\|^2}\). Substitute \(a=1/(4t)\) and \(c=x-tg_x\). The normalizers agree by definition, and \(-a\|y-c\|^2=-\|y-c\|^2/(4t)\). This pointwise identity makes the time and spatial calculations use precisely the same normalized Gaussian kernel.



Both kernels have the same Gaussian center and covariance. Substituting their normalizers and exponents proves equality pointwise, allowing spatial and time results developed in the two notations to be combined.

\begin{leanbox}{Lean: \texttt{frozenGaussianKernel\_\allowbreak{}eq\_\allowbreak{}realGaussianKernel}}
\leansource{\texttt{frozenGaussianKernel\_\allowbreak{}eq\_\allowbreak{}realGaussianKernel} \quad UnadjustedLangevin/\allowbreak{}LSI/\allowbreak{}GaussianTimeDerivative.lean, lines 279–284}
\begin{leancode}
\leanline{0}{theorem\ frozenGaussianKernel\_eq\_realGaussianKernel\ (U\ :\ E\ →\ ℝ)\ (t\ :\ ℝ)\ (x\ y\ :\ E)\ :}
\leanline{4}{frozenGaussianKernel\ U\ t\ x\ y\ =}
\leanline{6}{realGaussianKernel\ (1\ /\ (4\ *\ t))\ (x\ -\ t\ •\ gradient\ U\ x)\ y}
\end{leancode}
\leanlabel{How the Lean proof follows the argument}
\texttt{frozen\allowbreak{}Gaussian\allowbreak{}Kernel\_\allowbreak{}eq\_\allowbreak{}real\allowbreak{}Gaussian\allowbreak{}Kernel}: the proof is the application \texttt{unfold frozen\allowbreak{}Gaussian\allowbreak{}Kernel real\allowbreak{}Gaussian\allowbreak{}Kernel ULA.\allowbreak{}gaussian\allowbreak{}Proposal\allowbreak{}Normalizer}. It specializes that result with the hypotheses and objects in the exact statement.
\end{leanbox}


Earlier results used here: \cite[Section~3.37.12]{LangevinLean2026ULASC}.

\section{\texorpdfstring{A weighted derivative bound}{A weighted derivative bound}}


Retain the Gaussian factor instead of replacing it by its supremum. With \(r=y-x+tg_x\), the bound \(\|y-x\|^2\le2\|r\|^2+2t^2\|g_x\|^2\) implies that the absolute derivative coefficient is at most \(d/(2t)+\|r\|^2/(2t^2)+\|g_x\|^2\). Multiplication by \(K_t\ge0\) yields \(|\dot K_t|\le K_t[d/(2t)+\|r\|^2/(2t^2)+\|g_x\|^2]\). Unlike the initial-variable envelope, this bound remains integrable in the output variable.



The weighted bound retains the kernel's Gaussian decay and separates the initial-gradient square from displacement-dependent terms. This is the form suitable for integrating first in the observation.

\begin{leanbox}{Lean: \texttt{norm\_\allowbreak{}frozenGaussianKernelDot\_\allowbreak{}le\_\allowbreak{}weighted}}
\leansource{\texttt{norm\_\allowbreak{}frozenGaussianKernelDot\_\allowbreak{}le\_\allowbreak{}weighted} \quad UnadjustedLangevin/\allowbreak{}LSI/\allowbreak{}GaussianTimeDerivative.lean, lines 288–319}
\begin{leancode}
\leanline{0}{theorem\ norm\_frozenGaussianKernelDot\_le\_weighted}
\leanline{4}{(U\ :\ E\ →\ ℝ)\ (t\ :\ ℝ)\ (ht\ :\ 0\ <\ t)\ (x\ y\ :\ E)\ :}
\leanline{4}{‖frozenGaussianKernelDot\ U\ t\ x\ y‖\ ≤\ frozenGaussianKernel\ U\ t\ x\ y\ *}
\leanline{6}{((Module.finrank\ ℝ\ E\ :\ ℝ)\ /\ (2\ *\ t)\ +}
\leanline{8}{‖y\ -\ (x\ -\ t\ •\ gradient\ U\ x)‖\ \textasciicircum{}\ 2\ /\ (2\ *\ t\ \textasciicircum{}\ 2)\ +\ ‖gradient\ U\ x‖\ \textasciicircum{}\ 2)}
\end{leancode}
\leanlabel{How the Lean proof follows the argument}
\texttt{norm\_\allowbreak{}frozen\allowbreak{}Gaussian\allowbreak{}Kernel\allowbreak{}Dot\_\allowbreak{}le\_\allowbreak{}weighted}: the intermediate facts include \texttt{r}, \texttt{hvec}, \texttt{hn}, \texttt{hdis0}, \texttt{h\allowbreak{}K}, in source order; the source steps apply or rewrite \texttt{norm\_\allowbreak{}smul}, \texttt{Real.\allowbreak{}norm\_\allowbreak{}eq\_\allowbreak{}abs}, \texttt{norm\_\allowbreak{}mul}, \texttt{abs\_\allowbreak{}of\_\allowbreak{}nonneg}, \texttt{mul\_\allowbreak{}le\_\allowbreak{}mul\_\allowbreak{}of\_\allowbreak{}nonneg\_\allowbreak{}left}. The remaining scalar rearrangements are checked by \texttt{field\_\allowbreak{}simp}, \texttt{nlinarith}.
\end{leanbox}


Earlier results used here: \cite[Section~3.37.14]{LangevinLean2026ULASC}.

\section{\texorpdfstring{A Gaussian weight controlling the time derivative}{A Gaussian weight controlling the time derivative}}


Define the nonnegative weight \(\mathcal E_t(x,y)=K_t(x,y)(1+\|g_x\|^2+\|y-c_t(x)\|^2)\). Its three terms correspond to Gaussian mass, initial gradient energy, and Gaussian displacement energy. This weight dominates the derivative with a time-dependent scalar coefficient and can also be reused in spatial integration by parts.



The energy combines one, displacement square, and initial-gradient square under the Gaussian kernel. These are chosen to dominate every polynomial term appearing in the explicit time derivative.

\begin{leanbox}{Lean: \texttt{frozenGaussianEnergy}}
\leansource{\texttt{frozenGaussianEnergy} \quad UnadjustedLangevin/\allowbreak{}LSI/\allowbreak{}GaussianTimeDerivative.lean, lines 322–324}
\begin{leancode}
\leanline{0}{def\ frozenGaussianEnergy\ (U\ :\ E\ →\ ℝ)\ (t\ :\ ℝ)\ (x\ y\ :\ E)\ :\ ℝ\ :=}
\leanline{2}{frozenGaussianKernel\ U\ t\ x\ y\ *}
\leanline{4}{(1\ +\ ‖gradient\ U\ x‖\ \textasciicircum{}\ 2\ +\ ‖y\ -\ (x\ -\ t\ •\ gradient\ U\ x)‖\ \textasciicircum{}\ 2)}
\end{leancode}
\leanlabel{How the Lean proof follows the argument}
\texttt{frozen\allowbreak{}Gaussian\allowbreak{}Energy} is a definition with defining expression \texttt{frozen\allowbreak{}Gaussian\allowbreak{}Kernel U t x y *}. The right-hand side specifies the object used in the mathematical statement.
\end{leanbox}

\section{\texorpdfstring{Nonnegativity of the Gaussian weight}{Nonnegativity of the Gaussian weight}}


At positive time the normalizer and the Gaussian exponential are positive, so \(K_t\ge0\). Every term in \(1+\|g_x\|^2+\|y-c_t(x)\|^2\) is nonnegative. Their product is nonnegative; this identifies its norm with itself in the Fubini integrability proof.



Every component is a nonnegative square or constant multiplied by a nonnegative kernel. Their sum is therefore nonnegative, permitting nonnegative Fubini before finiteness is established.

\begin{leanbox}{Lean: \texttt{frozenGaussianEnergy\_\allowbreak{}nonneg}}
\leansource{\texttt{frozenGaussianEnergy\_\allowbreak{}nonneg} \quad UnadjustedLangevin/\allowbreak{}LSI/\allowbreak{}GaussianTimeDerivative.lean, lines 326–329}
\begin{leancode}
\leanline{0}{private\ theorem\ frozenGaussianEnergy\_nonneg\ (U\ :\ E\ →\ ℝ)\ (t\ :\ ℝ)\ (ht\ :\ 0\ <\ t)\ (x\ y\ :\ E)\ :}
\leanline{4}{0\ ≤\ frozenGaussianEnergy\ U\ t\ x\ y}
\end{leancode}
\leanlabel{How the Lean proof follows the argument}
\texttt{frozen\allowbreak{}Gaussian\allowbreak{}Energy\_\allowbreak{}nonneg}: the source steps apply or rewrite \texttt{mul\_\allowbreak{}nonneg}, \texttt{ULA.\allowbreak{}gaussian\allowbreak{}Proposal\allowbreak{}Normalizer\_\allowbreak{}pos}.
\end{leanbox}


Earlier results used here: \cite[Section~3.37.14]{LangevinLean2026ULASC}.

\section{\texorpdfstring{The component energy integral}{The component energy integral}}


For fixed \(x\), Gaussian normalization gives \(\int K_t(x,y)dy=1\), and the centered Gaussian covariance gives \(\int K_t(x,y)\|y-c_t(x)\|^2dy=2td\). Both integrands are integrable. The initial gradient factor is constant in \(y\), so its integral is \(\|g_x\|^2\). Expand the common energy into these three integrable terms to obtain \(\int\mathcal E_t(x,y)dy=1+\|g_x\|^2+2td\).



For each fixed initial state, integrate the kernel to one and its displacement square using the Gaussian moment formula. The resulting expression depends only on time, dimension, and the initial-gradient norm.

\begin{leanbox}{Lean: \texttt{frozenGaussianEnergy\_\allowbreak{}integral}}
\leansource{\texttt{frozenGaussianEnergy\_\allowbreak{}integral} \quad UnadjustedLangevin/\allowbreak{}LSI/\allowbreak{}GaussianTimeDerivative.lean, lines 332–359}
\begin{leancode}
\leanline{0}{theorem\ frozenGaussianEnergy\_integral\ (U\ :\ E\ →\ ℝ)\ (t\ :\ ℝ)\ (ht\ :\ 0\ <\ t)\ (x\ :\ E)\ :}
\leanline{4}{Integrable\ (frozenGaussianEnergy\ U\ t\ x)\ volume\ ∧}
\leanline{6}{(∫\ y,\ frozenGaussianEnergy\ U\ t\ x\ y)\ =}
\leanline{8}{1\ +\ ‖gradient\ U\ x‖\ \textasciicircum{}\ 2\ +\ 2\ *\ t\ *\ (Module.finrank\ ℝ\ E\ :\ ℝ)}
\end{leancode}
\leanlabel{How the Lean proof follows the argument}
\texttt{frozen\allowbreak{}Gaussian\allowbreak{}Energy\_\allowbreak{}integral}: the intermediate facts include \texttt{h\allowbreak{}K}, \texttt{h\allowbreak{}Ki}, \texttt{hmoment}, \texttt{hfun}, in source order; the source steps apply or rewrite \texttt{ULA.\allowbreak{}integrable\_\allowbreak{}radial\allowbreak{}Gaussian\allowbreak{}Density\_\allowbreak{}real}, \texttt{ULA.\allowbreak{}integral\_\allowbreak{}radial\allowbreak{}Gaussian\allowbreak{}Density\_\allowbreak{}real}, \texttt{hmoment.\allowbreak{}1}, \texttt{Pi.\allowbreak{}add\_\allowbreak{}def}, \texttt{integral\_\allowbreak{}add}. The remaining scalar rearrangements are checked by \texttt{ring}.
\end{leanbox}


Earlier results used here: \cite[Section~3.37.6]{LangevinLean2026ULASC}; \cite[Section~3.37.7]{LangevinLean2026ULASC}.

\section{\texorpdfstring{Integrability of the Gaussian weight in both variables}{Integrability of the Gaussian weight in both variables}}


The common energy is jointly measurable. For every initial \(x\), its output section is integrable by the component energy theorem. Since the energy is nonnegative, the integral of its norm is exactly \(1+\|g_x\|^2+2td\), an integrable function of \(x\) under the probability law \(\mu\). The product integrability criterion therefore proves \(\mathcal E_t\in L^1(\mu\otimes dy)\).



The initial-gradient square is integrable by the score-energy hypothesis, and all remaining terms are finite constants. The component integral calculation therefore proves integrability of the full pair energy.

\begin{leanbox}{Lean: \texttt{integrable\_\allowbreak{}frozenGaussianEnergy}}
\leansource{\texttt{integrable\_\allowbreak{}frozenGaussianEnergy} \quad UnadjustedLangevin/\allowbreak{}LSI/\allowbreak{}GaussianTimeDerivative.lean, lines 363–374}
\begin{leancode}
\leanline{0}{theorem\ integrable\_frozenGaussianEnergy}
\leanline{4}{(U\ :\ E\ →\ ℝ)\ (μ\ :\ Measure\ E)\ [IsProbabilityMeasure\ μ]}
\leanline{4}{(L\ :\ ℝ≥0)\ (hLip\ :\ LipschitzWith\ L\ (gradient\ U))}
\leanline{4}{(hG\ :\ Integrable\ (fun\ x\ ↦\ ‖gradient\ U\ x‖\ \textasciicircum{}\ 2)\ μ)\ (t\ :\ ℝ)\ (ht\ :\ 0\ <\ t)\ :}
\leanline{4}{Integrable\ (fun\ p\ :\ E\ ×\ E\ ↦\ frozenGaussianEnergy\ U\ t\ p.1\ p.2)\ (μ.prod\ volume)}
\end{leancode}
\leanlabel{How the Lean proof follows the argument}
\texttt{integrable\_\allowbreak{}frozen\allowbreak{}Gaussian\allowbreak{}Energy}: the intermediate fact is \texttt{hg}, in source order; the source steps apply or rewrite \texttt{integrable\_\allowbreak{}prod\_\allowbreak{}iff}, \texttt{fun\_\allowbreak{}prop}, \texttt{frozen\allowbreak{}Gaussian\allowbreak{}Energy\_\allowbreak{}integral}, \texttt{Pi.\allowbreak{}add\_\allowbreak{}def}, \texttt{integrable\_\allowbreak{}const}.
\end{leanbox}

\section{\texorpdfstring{The integral of the Gaussian weight}{The integral of the Gaussian weight}}


Joint integrability allows Fubini, giving \(\int\mathcal E_t\,d(\mu\otimes dy)=\int(1+\|g_x\|^2+2td)d\mu(x)\). Probability normalization and the definition of \(G_0\) reduce this to \(1+G_0+2td\). The displayed equality supplies a finite quantitative bound, in addition to proving integrability.



Integrating the component formula over the starting law replaces the squared gradient by \(G_0\). This gives an explicit finite energy expression, retaining its time dependence for uniform compact-interval estimates.

\begin{leanbox}{Lean: \texttt{integral\_\allowbreak{}frozenGaussianEnergy}}
\leansource{\texttt{integral\_\allowbreak{}frozenGaussianEnergy} \quad UnadjustedLangevin/\allowbreak{}LSI/\allowbreak{}GaussianTimeDerivative.lean, lines 377–387}
\begin{leancode}
\leanline{0}{theorem\ integral\_frozenGaussianEnergy}
\leanline{4}{(U\ :\ E\ →\ ℝ)\ (μ\ :\ Measure\ E)\ [IsProbabilityMeasure\ μ]}
\leanline{4}{(L\ :\ ℝ≥0)\ (hLip\ :\ LipschitzWith\ L\ (gradient\ U))}
\leanline{4}{(hG\ :\ Integrable\ (fun\ x\ ↦\ ‖gradient\ U\ x‖\ \textasciicircum{}\ 2)\ μ)\ (t\ :\ ℝ)\ (ht\ :\ 0\ <\ t)\ :}
\leanline{4}{(∫\ p\ :\ E\ ×\ E,\ frozenGaussianEnergy\ U\ t\ p.1\ p.2\ ∂μ.prod\ volume)\ =}
\leanline{6}{1\ +\ (∫\ x,\ ‖gradient\ U\ x‖\ \textasciicircum{}\ 2\ ∂μ)\ +\ 2\ *\ t\ *\ (Module.finrank\ ℝ\ E\ :\ ℝ)}
\end{leancode}
\leanlabel{How the Lean proof follows the argument}
\texttt{integral\_\allowbreak{}frozen\allowbreak{}Gaussian\allowbreak{}Energy}: the source steps apply or rewrite \texttt{integral\_\allowbreak{}prod}, \texttt{integrable\_\allowbreak{}frozen\allowbreak{}Gaussian\allowbreak{}Energy}, \texttt{integral\_\allowbreak{}add}, \texttt{Pi.\allowbreak{}add\_\allowbreak{}def}, \texttt{integrable\_\allowbreak{}const}.
\end{leanbox}

\section{\texorpdfstring{A uniform coefficient on a positive time interval}{A uniform coefficient on a positive time interval}}


Set \(C_t=d/(2t)+1/(2t^2)+1\). It is positive for \(t>0\), and is bounded above on any interval \([a,b]\) with \(a>0\) by \(C_a\). This finite coefficient suffices for the local space-time integrability required by the entropy chain rule.



On a positive compact interval, \(t\) is bounded above and \(1/t\) is bounded above by the reciprocal of the lower endpoint. One coefficient can therefore dominate all scalar time factors uniformly.

\begin{leanbox}{Lean: \texttt{gaussianTimeCoefficient}}
\leansource{\texttt{gaussianTimeCoefficient} \quad UnadjustedLangevin/\allowbreak{}LSI/\allowbreak{}GaussianTimeDerivative.lean, lines 391–392}
\begin{leancode}
\leanline{0}{def\ gaussianTimeCoefficient\ (t\ :\ ℝ)\ :\ ℝ\ :=}
\leanline{2}{(Module.finrank\ ℝ\ E\ :\ ℝ)\ /\ (2\ *\ t)\ +\ 1\ /\ (2\ *\ t\ \textasciicircum{}\ 2)\ +\ 1}
\end{leancode}
\leanlabel{How the Lean proof follows the argument}
\texttt{gaussian\allowbreak{}Time\allowbreak{}Coefficient} is a definition with defining expression \texttt{(Module.\allowbreak{}finrank \ensuremath{\mathbb R} E : \ensuremath{\mathbb R}) /\allowbreak{} (2 * t) + 1 /\allowbreak{} (2 * t \textasciicircum{} 2) + 1}. The right-hand side specifies the object used in the mathematical statement.
\end{leanbox}

\section{\texorpdfstring{Derivative domination by the Gaussian weight}{Derivative domination by the Gaussian weight}}


In the weighted derivative bound, denote the three nonnegative coefficients by \(A=d/(2t)\), \(B=1/(2t^2)\), and \(1\). Each is at most \(C_t=A+B+1\). Therefore \(A+B\|r\|^2+\|g_x\|^2\le C_t(1+\|g_x\|^2+\|r\|^2)\). Multiply by \(K_t\ge0\) to prove \(|\dot K_t|\le C_t\mathcal E_t\).



Multiply that coefficient by the nonnegative Gaussian weight. The pointwise derivative bound then holds throughout the interval with one common integrable function, which is the domination hypothesis for the space-time integral.

\begin{leanbox}{Lean: \texttt{norm\_\allowbreak{}frozenGaussianKernelDot\_\allowbreak{}le\_\allowbreak{}energy}}
\leansource{\texttt{norm\_\allowbreak{}frozenGaussianKernelDot\_\allowbreak{}le\_\allowbreak{}energy} \quad UnadjustedLangevin/\allowbreak{}LSI/\allowbreak{}GaussianTimeDerivative.lean, lines 395–425}
\begin{leancode}
\leanline{0}{theorem\ norm\_frozenGaussianKernelDot\_le\_energy\ (U\ :\ E\ →\ ℝ)\ (t\ :\ ℝ)}
\leanline{4}{(ht\ :\ 0\ <\ t)\ (x\ y\ :\ E)\ :}
\leanline{4}{‖frozenGaussianKernelDot\ U\ t\ x\ y‖\ ≤}
\leanline{6}{gaussianTimeCoefficient\ (E\ :=\ E)\ t\ *\ frozenGaussianEnergy\ U\ t\ x\ y}
\end{leancode}
\leanlabel{How the Lean proof follows the argument}
\texttt{norm\_\allowbreak{}frozen\allowbreak{}Gaussian\allowbreak{}Kernel\allowbreak{}Dot\_\allowbreak{}le\_\allowbreak{}energy}: the intermediate facts include \texttt{h}, \texttt{A}, \texttt{B}, \texttt{h\ensuremath{{}_{2}}}, \texttt{h\ensuremath{{}_{3}}}, in source order; the source steps apply or rewrite \texttt{add\_\allowbreak{}le\_\allowbreak{}add}, \texttt{mul\_\allowbreak{}le\_\allowbreak{}mul\_\allowbreak{}of\_\allowbreak{}nonneg\_\allowbreak{}right}, \texttt{one\_\allowbreak{}mul}. The \texttt{calc} block checks the successive equalities or inequalities in the argument. The remaining scalar rearrangements are checked by \texttt{ring}, \texttt{linarith}.
\end{leanbox}


Earlier results used here: \cite[Section~3.37.14]{LangevinLean2026ULASC}.

\section{\texorpdfstring{Integrability of the kernel time derivative}{Integrability of the kernel time derivative}}


The component time derivative is jointly measurable, and its norm is bounded by \(C_t\mathcal E_t\). The common energy is integrable on \(\mu\otimes dy\), and \(C_t\) is a finite scalar at positive time. Domination therefore proves joint integrability of \(\dot K_t\), which licenses the changes of integration order used in the time weak equation.



The energy envelope has finite pair integral and the derivative is measurable. Absolute-value domination consequently proves integrability of the kernel time derivative under the Gaussian pair measure.

\begin{leanbox}{Lean: \texttt{integrable\_\allowbreak{}frozenGaussianKernelDot}}
\leansource{\texttt{integrable\_\allowbreak{}frozenGaussianKernelDot} \quad UnadjustedLangevin/\allowbreak{}LSI/\allowbreak{}GaussianTimeDerivative.lean, lines 429–439}
\begin{leancode}
\leanline{0}{theorem\ integrable\_frozenGaussianKernelDot}
\leanline{4}{(U\ :\ E\ →\ ℝ)\ (μ\ :\ Measure\ E)\ [IsProbabilityMeasure\ μ]}
\leanline{4}{(L\ :\ ℝ≥0)\ (hLip\ :\ LipschitzWith\ L\ (gradient\ U))}
\leanline{4}{(hG\ :\ Integrable\ (fun\ x\ ↦\ ‖gradient\ U\ x‖\ \textasciicircum{}\ 2)\ μ)\ (t\ :\ ℝ)\ (ht\ :\ 0\ <\ t)\ :}
\leanline{4}{Integrable\ (fun\ p\ :\ E\ ×\ E\ ↦\ frozenGaussianKernelDot\ U\ t\ p.1\ p.2)\ (μ.prod\ volume)}
\end{leancode}
\leanlabel{How the Lean proof follows the argument}
\texttt{integrable\_\allowbreak{}frozen\allowbreak{}Gaussian\allowbreak{}Kernel\allowbreak{}Dot}: the intermediate fact is \texttt{hg}, in source order; the source steps apply or rewrite \texttt{integrable\_\allowbreak{}frozen\allowbreak{}Gaussian\allowbreak{}Energy}, \texttt{const\_\allowbreak{}mul}, \texttt{ae\_\allowbreak{}of\_\allowbreak{}all}, \texttt{norm\_\allowbreak{}frozen\allowbreak{}Gaussian\allowbreak{}Kernel\allowbreak{}Dot\_\allowbreak{}le\_\allowbreak{}energy}.
\end{leanbox}

\section{\texorpdfstring{Spatial integrability of the mixture derivative}{Spatial integrability of the mixture derivative}}


Fubini applied to the jointly integrable component derivative first gives \(\dot p_t\in L^1(dy)\). The norm inequality for Bochner integration gives \(|\dot p_t(y)|\le\int|\dot K_t(x,y)|d\mu(x)\). Integrate in \(y\), exchange the two integrals, and dominate by \(C_t\mathcal E_t\). The exact envelope integral yields \(\int|\dot p_t(y)|dy\le C_t(1+G_0+2td)\).



The norm of the mixture derivative is at most the mixture of the kernel derivative norms. Fubini and the finite pair bound therefore give spatial integrability after the initial state is averaged out.

\begin{leanbox}{Lean: \texttt{integrable\_\allowbreak{}interpolationLebesgueDensityTimeDerivative}}
\leansource{\texttt{integrable\_\allowbreak{}interpolationLebesgueDensityTimeDerivative} \quad UnadjustedLangevin/\allowbreak{}LSI/\allowbreak{}GaussianTimeDerivative.lean, lines 444–471}
\begin{leancode}
\leanline{0}{theorem\ integrable\_interpolationLebesgueDensityTimeDerivative}
\leanline{4}{(U\ :\ E\ →\ ℝ)\ (μ\ :\ Measure\ E)\ [IsProbabilityMeasure\ μ]}
\leanline{4}{(L\ :\ ℝ≥0)\ (hLip\ :\ LipschitzWith\ L\ (gradient\ U))}
\leanline{4}{(hG\ :\ Integrable\ (fun\ x\ ↦\ ‖gradient\ U\ x‖\ \textasciicircum{}\ 2)\ μ)\ (t\ :\ ℝ)\ (ht\ :\ 0\ <\ t)\ :}
\leanline{4}{Integrable\ (interpolationLebesgueDensityTimeDerivative\ U\ μ\ t)\ volume\ ∧}
\leanline{6}{(∫\ y,\ ‖interpolationLebesgueDensityTimeDerivative\ U\ μ\ t\ y‖)\ ≤}
\leanline{8}{gaussianTimeCoefficient\ (E\ :=\ E)\ t\ *}
\leanline{10}{(1\ +\ (∫\ x,\ ‖gradient\ U\ x‖\ \textasciicircum{}\ 2\ ∂μ)\ +\ 2\ *\ t\ *\ (Module.finrank\ ℝ\ E\ :\ ℝ))}
\end{leancode}
\leanlabel{How the Lean proof follows the argument}
\texttt{integrable\_\allowbreak{}interpolation\allowbreak{}Lebesgue\allowbreak{}Density\allowbreak{}Time\allowbreak{}Derivative}: the intermediate facts include \texttt{hprod}, \texttt{hi}, in source order; the source steps apply or rewrite \texttt{integral\_\allowbreak{}mono}, \texttt{hi.\allowbreak{}norm}, \texttt{hprod.\allowbreak{}norm}, \texttt{integral\_\allowbreak{}const\_\allowbreak{}mul}, \texttt{integral\_\allowbreak{}frozen\allowbreak{}Gaussian\allowbreak{}Energy}. The \texttt{calc} block checks the successive equalities or inequalities in the argument.
\end{leanbox}

\section{\texorpdfstring{Spacetime integrability of the mixture derivative}{Spacetime integrability of the mixture derivative}}


Joint measurability of \(\dot p\) permits the product integrability criterion on \([a,b]\times E\), where \(a>0\). Every time section is in \(L^1(dy)\). For times in the interval, \(C_t\le C_a\) and \(1+G_0+2td\le1+G_0+2bd\), giving the common finite scalar bound \(\int|\dot p_t|dy\le C_a(1+G_0+2bd)\). The time interval has finite measure, so the section norms are integrable. Hence \(\dot p\in L^1(dt\,dy)\) locally away from zero. If \(a>b\), the restricted time measure is zero and the same argument remains valid.



The uniform compact-time coefficient and the energy integral provide an integrable time majorant. Integrating that majorant over the finite interval proves the full spacetime \(L^1\) estimate used by bounded entropy calculus.

\begin{leanbox}{Lean: \texttt{integrable\_\allowbreak{}interpolationLebesgueDensityTimeDerivative\_\allowbreak{}spacetime}}
\leansource{\texttt{integrable\_\allowbreak{}interpolationLebesgueDensityTimeDerivative\_\allowbreak{}spacetime} \quad UnadjustedLangevin/\allowbreak{}LSI/\allowbreak{}GaussianTimeDerivative.lean, lines 476–505}
\begin{leancode}
\leanline{0}{theorem\ integrable\_interpolationLebesgueDensityTimeDerivative\_spacetime}
\leanline{4}{(U\ :\ E\ →\ ℝ)\ (μ\ :\ Measure\ E)\ [IsProbabilityMeasure\ μ]}
\leanline{4}{(L\ :\ ℝ≥0)\ (hLip\ :\ LipschitzWith\ L\ (gradient\ U))}
\leanline{4}{(hG\ :\ Integrable\ (fun\ x\ ↦\ ‖gradient\ U\ x‖\ \textasciicircum{}\ 2)\ μ)}
\leanline{4}{(a\ b\ :\ ℝ)\ (ha\ :\ 0\ <\ a)\ :}
\leanline{4}{Integrable\ (fun\ p\ :\ ℝ\ ×\ E\ ↦\ interpolationLebesgueDensityTimeDerivative\ U\ μ\ p.1\ p.2)}
\leanline{6}{((volume.restrict\ (Icc\ a\ b)).prod\ volume)}
\end{leancode}
\leanlabel{How the Lean proof follows the argument}
\texttt{integrable\_\allowbreak{}interpolation\allowbreak{}Lebesgue\allowbreak{}Density\allowbreak{}Time\allowbreak{}Derivative\_\allowbreak{}spacetime}: the intermediate facts include \texttt{C}, \texttt{hm}, \texttt{htpos}, \texttt{hd}, \texttt{hmom}, in source order; the source steps apply or rewrite \texttt{integrable\_\allowbreak{}prod\_\allowbreak{}iff}, \texttt{hm.\allowbreak{}aestrongly\allowbreak{}Measurable}, \texttt{norm\_\allowbreak{}nonneg}, \texttt{hbound.\allowbreak{}trans}, \texttt{mul\_\allowbreak{}le\_\allowbreak{}mul}.
\end{leanbox}

\begin{thebibliography}{ULA-LSI}
\bibitem[ULA-SC]{LangevinLean2026ULASC}
LangevinLean Project. \emph{Unadjusted Langevin under Strong Convexity}. LL-01, 2026. Draft report, content revision d34ead06987c. \url{https://odakyildiz.github.io/LangevinLean/docs/pdf/UnadjustedLangevin/FullProof.pdf}.
\end{thebibliography}

\end{document}
